Hybrid and all-inorganic halide perovskites have emerged as leading
candidates for next-generation photovoltaic (PV) absorbers, with
certified single-junction power-conversion efficiencies now
exceeding \SI{26}{\percent}.\cite{ref_efficiency}
A persistent challenge in translating laboratory performance to
commercial deployment is the toxicity of lead, which remains the
standard B-site cation in high-efficiency perovskite absorbers.
\cite{ref_pb_toxicity}
Tin-based perovskites --- including the all-inorganic \CsSnI{} and
the hybrid methylammonium (MA) and formamidinium (FA) variants
\MASnI{} and \FASnI{} --- have attracted considerable attention as
lead-free alternatives, owing to their favourable optical band gaps
in the \SIrange{1.2}{1.4}{eV} range, high charge-carrier mobilities,
and structural versatility.\cite{pitaro2022tin, stoumpos2013semiconducting}
However, tin-halide perovskites exhibit pronounced chemical
instability under ambient conditions, strong polymorphism across
accessible temperature ranges, and a complex interplay between
structure and electronic properties that complicates both materials
optimisation and device integration.\cite{ref_sn_instability}

Among the properties most directly relevant to PV performance, the
optical band gap is central.
In a solar-cell absorber, the band gap governs both the photon
absorption onset and the thermodynamic ceiling on open-circuit
voltage, making its temperature dependence a practically important
quantity.
Real devices experience a wide temperature range during fabrication,
module lamination, and outdoor service, and processing steps such as
annealing can transiently expose the absorber layer to temperatures
at which structural phase changes occur.\cite{ref_processing_temp}
Understanding how the band gap responds to temperature --- and which
structural degrees of freedom drive that response --- is therefore
not merely a fundamental question but one with direct manufacturing
relevance.

Temperature affects the electronic structure of halide perovskites
through at least three coupled mechanisms.
First, thermal expansion dilates the unit cell, which modifies
metal--halide bond lengths and bond angles and directly perturbs
the band edges.\cite{huang2013electronic}
Second, tin-halide perovskites are strongly polymorphic:
\CsSnI{}, \MASnI{}, and \FASnI{} each undergo transitions among
orthorhombic, tetragonal, and cubic phases across the temperature
range relevant to PV processing and operation, and each polymorph
carries a distinct pattern of octahedral tilting and connectivity
that shifts the electronic structure independently of lattice
volume.\cite{schueller2018crystal, stoumpos2013semiconducting}
Third, in hybrid perovskites, the organic A-site cations --- MA and
FA --- are rotationally mobile, and their orientational disorder
couples to the inorganic \ce{Sn-I} cage, introducing an additional
structural degree of freedom that can modulate the band gap through
A-site--cage interactions.\cite{laven2023rotational, maheshwari2019relation}
In practice, these three contributions act simultaneously, making it
difficult to assign observed temperature-dependent band-gap changes
to a single mechanism.

Prior computational studies have individually addressed
lattice-constant effects on band gap,\cite{huang2013electronic}
phase-dependent electronic structure in selected tin
perovskites,\cite{stoumpos2013semiconducting} and cation rotational
dynamics in hybrid systems.\cite{maheshwari2019relation, hautzinger2024asite}
However, a systematic study that separates and quantifies all three
contributions --- expansion, phase identity, and cation disorder ---
within a common computational framework, and across materials
spanning both all-inorganic and hybrid chemistries, has not been
reported.
The absence of such a study leaves open questions that are directly
relevant to materials design: does lattice expansion dominate over
phase effects, or vice versa?
Do MA and FA perturb the band gap in the same way or through
distinct mechanisms?
And how do these effects combine under realistic finite-temperature
conditions that include simultaneous lattice dilation and structural
disorder?

In this work, we address these questions through a controlled
computational study of \CsSnI{}, \MASnI{}, and \FASnI{} across
their three accessible polymorphs.
We deploy a validated MLFF for long-time finite-temperature MD to
characterise phase behaviour, and a system-specific accelerated
single-step DFT (SS-DFT) framework for efficient band-gap evaluation
over large supercells.
The study is organised around four controlled tasks:
\begin{enumerate}
  \item finite-temperature phase-transition analysis via long-time
        MD heating and cooling cycles;
  \item controlled lattice-scaling scans to isolate the thermal
        expansion contribution to band gap;
  \item cation-orientation sensitivity studies in the two hybrid
        systems; and
  \item an integrated demonstration combining finite-temperature
        structural sampling with externally imposed strain, intended
        to bridge the gap between idealised mechanism studies and
        more realistic process-relevant conditions.
\end{enumerate}

Our results demonstrate that the net temperature dependence of band
gap in these tin perovskites is a composite effect in which lattice
expansion, phase identity, and cation orientation contribute
independently and can either reinforce or partially cancel one another
depending on the composition and temperature regime.
These findings provide a mechanistic foundation for rational
band-gap engineering in lead-free tin-halide perovskites and
establish a reusable computational workflow applicable to a broader
class of halide perovskite systems.
